% ============================================================
%  SECCIÓN SANTA FE — Venta de boletería Independiente Santa Fe
%  Generada automáticamente. Listo para pegar en el paper.
% ============================================================

\subsection{Descripción del escenario: Boletería Independiente Santa Fe}

Consideremos el sistema de venta de boletas para los partidos del
\textbf{Independiente Santa Fe} en el estadio Nemesio Camacho El Campín.
En este contexto, cada taquilla física del estadio opera como un servidor
que atiende a un hincha a la vez. El sistema no posee fila de espera:
si un hincha llega y todas las taquillas están ocupadas, simplemente
se da la vuelta y se va sin comprar boleto.

La analogía formal con el modelo M/M/k/k se establece en la
Tabla~\ref{tab:analogia_santafe}.

\begin{table}[H]
\centering
\caption{Correspondencia entre boletería Santa Fe y el modelo M/M/k/k.}
\label{tab:analogia_santafe}
\begin{tabular}{ll}
\toprule
\textbf{Elemento del modelo M/M/k/k} & \textbf{Elemento en boletería Santa Fe} \\
\midrule
Tasa de llegadas $\lambda$ & Tasa de llegada de hinchas (fans/min) \\
Tiempo de servicio $\sim \mathrm{Exp}(\mu)$ & Tiempo de venta por taquilla (ventas/min) \\
$k$ servidores en paralelo & $k$ taquillas abiertas en el estadio \\
Capacidad $=k$ (sin cola) & Sin fila de espera habilitada \\
Estado $n$: ocupación del sistema & Número de taquillas atendiendo hinchas \\
Bloqueo (prob.\ $B(k,a)$) & Proporción de hinchas que se van sin boleto \\
Carga ofrecida $a=\lambda/\mu$ & Intensidad de tráfico en Erlangs \\
\bottomrule
\end{tabular}
\end{table}

\subsection{Escenarios de partido}

Se analizan cuatro escenarios representativos del calendario del club,
ordenados por intensidad de demanda:

\begin{enumerate}
    \item \textbf{Partido de Liga regular:} partido contra rival de media tabla,
          con afluencia moderada de público.
    \item \textbf{Clásico vs Millonarios:} el partido más importante del año,
          con demanda extremadamente alta.
    \item \textbf{Final de Copa:} partido de eliminación directa, con afluencia
          masiva de hinchas.
    \item \textbf{Súper-clásico (Clausura):} definición del torneo contra el
          máximo rival, con demanda histórica.
\end{enumerate}

\subsection{Resultados de simulación}

La Tabla~\ref{tab:santa_fe_resultados} resume los parámetros de cada escenario
y la probabilidad de bloqueo obtenida por simulación de eventos discretos,
comparada con la predicción teórica de Erlang~B.

\begin{table}[H]
\centering
\caption{Escenarios de boletería Santa Fe y resultados de simulación M/M/k/k.}
\label{tab:santa_fe_resultados}
\begin{tabular}{lccccc}
\toprule
\textbf{Escenario} & $k$ & $\lambda$ & $\mu$ & $a$ & $\hat{B}_{\mathrm{sim}}$ \\
\midrule
Partido de Liga regular & 5 & 20 & 4.0 & 5.00 & 0.2855 \\
Clásico vs Millonarios & 10 & 60 & 3.0 & 20.00 & 0.5384 \\
Final de Copa & 15 & 100 & 2.0 & 50.00 & 0.7059 \\
Súper-clásico (Clausura) & 20 & 150 & 2.5 & 60.00 & 0.6748 \\
\bottomrule
\end{tabular}
\end{table}

\begin{remark}
La concordancia entre el valor teórico $B(k,a)$ y la simulación $\hat{B}_{\mathrm{sim}}$
es prácticamente exacta en todos los escenarios (diferencia menor al $0.01$),
lo que valida tanto la implementación del simulador como la aplicabilidad del
modelo M/M/k/k al caso de boletería.
\end{remark}

\subsection{Comparación gráfica: teoría vs simulación}

La Figura~\ref{fig:santafe_comparacion} muestra la probabilidad de bloqueo
para cada escenario, comparando la predicción de Erlang~B con el resultado
empírico de la simulación. La coincidencia es prácticamente perfecta.

\begin{figure}[H]
\centering
\includegraphics[width=0.95\textwidth]{plots/act_6_santafe/comparacion_bloqueo.png}
\caption{Comparación entre Erlang B analítico y simulación por eventos discretos
para los escenarios de boletería Santa Fe.}
\label{fig:santafe_comparacion}
\end{figure}

\subsection{Curvas de sensibilidad $B(k,a)$ vs $\lambda$}

La Figura~\ref{fig:santafe_curvas} ilustra cómo crece la probabilidad de
bloqueo al aumentar la tasa de llegada $\lambda$, manteniendo fijos $k$ y $\mu$
en cada escenario. Se destaca el punto de operación actual y la línea de
referencia $B_{\max}=0.05$ (umbral estándar de calidad de servicio).

\begin{figure}[H]
\centering
\includegraphics[width=0.98\textwidth]{plots/act_6_santafe/curvas_erlang_por_escenario.png}
\caption{Curvas de Erlang B para cada escenario de boletería. El punto negro
indica la operación actual; la línea punteada roja marca el umbral del $5\%$.}
\label{fig:santafe_curvas}
\end{figure}

\subsection{Evolución temporal de la ocupación}

La Figura~\ref{fig:santafe_ocupacion} muestra la evolución temporal del número
de taquillas ocupadas en una ventana de 30 minutos, para cada escenario.
Se observa claramente cómo en los escenarios de alta demanda (Final y
Súper-clásico) las taquillas permanecen casi permanentemente saturadas.

\begin{figure}[H]
\centering
\includegraphics[width=0.98\textwidth]{plots/act_6_santafe/ocupacion_temporal_todos.png}
\caption{Evolución temporal de la ocupación de taquillas en los primeros 30 minutos
de venta para cada escenario. La línea roja punteada indica la capacidad máxima $k$.}
\label{fig:santafe_ocupacion}
\end{figure}

\subsection{Dimensionamiento óptimo de taquillas}

La Tabla~\ref{tab:santa_fe_dimensionamiento} presenta el número mínimo de
taquillas $k^*$ necesarias para cumplir distintos niveles de calidad de servicio,
determinado mediante la recursión de Jagerman~(\ref{eq:jagerman}).

\begin{table}[H]
\centering
\caption{Dimensionamiento óptimo de taquillas por nivel de servicio.}
\label{tab:santa_fe_dimensionamiento}
\begin{tabular}{l|cccc}
\toprule
& \multicolumn{4}{c}{\textbf{$k^*$ mínimo para umbral $B_{\max}$}} \\
\textbf{Escenario} & $0.001$ & $0.01$ & $0.05$ & $0.10$ \\
\midrule
Partido de Liga regular & 14 & 11 & 9 & 8 \\
Clásico vs Millonarios & 35 & 30 & 26 & 23 \\
Final de Copa & 71 & 64 & 56 & 51 \\
Súper-clásico (Clausura) & 83 & 75 & 66 & 60 \\
\bottomrule
\end{tabular}
\end{table}

\begin{figure}[H]
\centering
\includegraphics[width=0.95\textwidth]{plots/act_6_santafe/dimensionamiento_optimo.png}
\caption{Dimensionamiento óptimo de taquillas por nivel de servicio.
El rombo rojo indica la configuración actual de cada escenario.}
\label{fig:santafe_dimensionamiento}
\end{figure}

\begin{remark}
Los resultados revelan que la configuración actual de taquillas es
\emph{insuficiente} para todos los escenarios analizados. Incluso para el
umbral más relajado ($B_{\max}=0.10$), el Partido de Liga regular requiere
$k^*=8$ taquillas (frente a las $k=5$ actuales), mientras que el Súper-clásico
necesitaría $k^*=60$ (frente a las $k=20$ actuales). En otras palabras:
\textbf{hay que ponerse las pilas con la logística}, hermano.
\end{remark}

\subsection{Distribución estacionaria $P_n$}

La Figura~\ref{fig:santafe_distribucion} muestra la distribución estacionaria
de probabilidades $P_n$ para cada escenario. La barra roja destaca $P_k$,
la probabilidad de que todas las taquillas estén ocupadas (bloqueo).

\begin{figure}[H]
\centering
\includegraphics[width=0.98\textwidth]{plots/act_6_santafe/distribucion_estacionaria.png}
\caption{Distribución estacionaria $P_n$ para cada escenario de boletería.
La barra roja indica $P_k$ (probabilidad de bloqueo).}
\label{fig:santafe_distribucion}
\end{figure}

\subsection{Código MATLAB}

La implementación completa de la simulación, incluyendo la generación de
datos, la comparación teoría vs simulación, las curvas de sensibilidad y el
dimensionamiento óptimo, se encuentra en el archivo
\texttt{matlab/actividad6\_santafe.m}. El script sigue la misma metodología
que las Actividades~3 y~6 del curso, adaptada al contexto de boletería.

% ============================================================
% FIN SECCIÓN SANTA FE
% ============================================================
